\chapter{Story Summary}
\label{app:story-summary}

\section{Story Points}
\label{app:story-summary:points}

\paragraph{Q1. What is the central question?}
How can conditional diffusion models be effectively leveraged as generative classifiers for out-of-distribution detection, and what design choices, particularly regarding the training objective (separation loss) and inference strategy (Monte Carlo timestep sampling, contrastive scoring), most strongly influence detection performance?

\paragraph{Q2. Why is this question important?}
Deployed deep learning models encounter inputs outside their training distribution yet produce overconfident predictions, creating dangerous failure modes in safety-critical applications such as healthcare, autonomous driving, and manufacturing. Existing generative OOD detectors suffer from the likelihood paradox, where models assign higher likelihood to OOD data than to their own training data. Diffusion models offer reconstruction-based scoring that sidesteps this failure mode, but systematic investigation of their design space for OOD detection is lacking.

\paragraph{Q3. What evidence/data (variables) are needed to answer this question?}
We need OOD detection metrics (AUROC, FPR@95\%TPR, AUPRC) across diverse near- and far-OOD benchmarks (CIFAR-100, Places365, FashionMNIST, SVHN, Textures) using a binary conditional diffusion model on CIFAR-10. Additionally, we require ablation results for the separation loss coefficient $\lambda$, the number of Monte Carlo trials $K$, and the scoring strategy (contrastive vs.\ ID-only). For industrial validation, we need public quality-classification AUROC, per-feature results, and separation-loss ablations from the YOLO+CDM inkjet pipeline under 5-fold cross-validation, all reproducible from \texttt{InkjetOOD}, public FTI\_Zer0P data, and released CDM weights.


\paragraph{Q4. What methods are used to get this evidence/data?}
We train a binary conditional diffusion model (DDPM with UNet backbone) on CIFAR-10 with class-conditional noise prediction and a novel separation loss that maximises the distance between class-conditional noise predictions, indirectly widening reconstruction-error gaps. OOD scoring uses Monte Carlo timestep sampling with contrastive difference scoring. For the industrial track, we build a public two-stage YOLO + CDM pipeline with multi-head conditioning and evaluate it using rigorous 5-fold cross-validation with image-level splitting.

\paragraph{Q5. What analyses must be applied to the data to answer the central question?}
Systematic ablation studies isolate the effect of each design choice: a $\lambda$-sweep ($\{0.0, 0.001, 0.01, 0.02, 0.05, 0.1\}$) quantifies the separation loss contribution; a $K$-sweep ($\{1, 5, 10, 25, 50, 100\}$) characterises the accuracy--efficiency frontier; and a scoring strategy comparison (ID-only vs.\ difference vs.\ ratio) establishes the necessity of contrastive scoring. Cross-domain analysis compares the separation loss effect between CIFAR-10 and inkjet data to identify boundary conditions. Statistical testing (paired $t$-tests with Holm correction for $n \geq 5$ folds) validates industrial results.

\paragraph{Q6. What evidence/data (values for the variables) were obtained?}
On CIFAR-10, the binary CDM achieves 98.98\% AUROC on the within-CIFAR split and 90.50--96.97\% across five external OOD datasets. The separation loss improves the three-seed Within-CIFAR mean by $+6.5$ pp (from $92.52\% \pm 11.07\%$ at $\lambda{=}0$ to $99.03\% \pm 0.07\%$ at $\lambda{=}0.02$) while substantially reducing seed sensitivity. On inkjet data, the public YOLO+CDM pipeline reaches $0.8673 \pm 0.0230$ AUROC at $\lambda{=}0$ under 5-fold cross-validation, with no statistically significant gain from non-zero separation-loss settings. Per-feature analysis shows that distance and dot features are easiest, while edge roughness features remain more difficult and variable.

\paragraph{Q7. What were the results of the analyses?}
For CIFAR-10 OOD detection, the separation loss is the clearest performance and stability driver: binary conditioning alone can work well but is highly seed-sensitive, while $\lambda{=}0.02$ raises the three-seed mean to $99.03\% \pm 0.07\%$. Contrastive scoring is essential, since using ID error alone collapses to 20.2\% AUROC on SVHN. $K{=}10$ Monte Carlo trials provide the best accuracy--efficiency trade-off (98.2\% AUROC at $5\times$ speedup over $K{=}50$). On inkjet data, the separation loss shows no statistically significant effect (Holm-adjusted $p > 0.05$), revealing that its benefit requires sufficient dataset scale and inter-class visual difference.

\paragraph{Q8. How did the analyses answer the central question?}
Conditional diffusion models are effective generative classifiers for OOD detection when three design choices are made correctly in the evaluated CIFAR-10 setting: (1) the training objective benefits from a separation loss that widens class-conditional reconstruction-error gaps; (2) contrastive scoring comparing reconstruction under both conditions is essential; and (3) Monte Carlo timestep averaging with $K \geq 10$ is needed for reliable performance. The industrial application confirms that reconstruction error generalises as a quality signal, but also shows that the separation-loss gain does not transfer automatically to smaller domain-specific datasets. On the inkjet industrial track, the public evidence supports YOLO+CDM performance, feature-level heterogeneity, and the negative separation-loss result, while broader method-ranking claims require complete public artefacts before being retained.

\paragraph{Q9. What does this answer tell us about the broader field?}
For tasks centred on distributional fit, such as OOD detection and quality classification, generative approaches modelling $p(x|y)$ through diffusion processes can provide strong, interpretable signals. The CIFAR-10 results suggest that objective-function engineering can matter as much as architecture when learning distributional boundaries. The negative separation-loss result on public inkjet data highlights that techniques validated on large-scale benchmarks do not automatically transfer to small domain-specific datasets, underscoring the importance of cross-domain validation.
